% !TEX root = ../matan.tex

\section{Дифференцирование сложной функции. Производная по направлению, градиент}

\subsection*{Дифференцирование сложной функции}

Напомним определение дифференцируемости. Пусть $\Omega\subset\RR^d$ открыто,
$x\in\Omega$, $f\colon\Omega\to\RR^n$. Функция $f$ \emph{дифференцируема} в точке
$x$, если существует линейный оператор $d_xf\colon\RR^d\to\RR^n$ такой, что при
$\xi\to x$, $\xi\in\Omega$,
\[
f(\xi)=f(x)+d_xf(\xi-x)+o(|\xi-x|).
\]
Оператор $d_xf$ называется \emph{дифференциалом} $f$ в точке $x$.

Рассмотрим композицию двух отображений
\[
f\colon\Omega\subset\RR^d\to\RR^n,
\qquad
g\colon\widetilde\Omega\subset\RR^n\to\RR^m,
\qquad
f(\Omega)\subset\widetilde\Omega,
\]
где $x$ --- внутренняя точка $\Omega$, $y=f(x)$ --- внутренняя точка $\widetilde\Omega$.

\begin{Theorem}{Производная сложной функции (правило цепочки)}{}
	Если $f$ дифференцируема в точке $x$, а $g$ дифференцируема в точке $y=f(x)$,
	то композиция $h=g\circ f$ дифференцируема в точке $x$ и
	\[
	d_x h = d_{f(x)}g\circ d_xf.
	\]
	При этом $d_xf\colon\RR^d\to\RR^n$, $d_yg\colon\RR^n\to\RR^m$, так что
	$d_x h\colon\RR^d\to\RR^m$.
\end{Theorem}

\begin{proof}
	Обозначим $A_x:=d_xf$ и $B_y:=d_yg$. Нужно показать, что
	\[
	h(\xi)-h(x)=(B_y\circ A_x)(\xi-x)+o(|\xi-x|).
	\]

	\textit{Шаг 1 (дифференцируемость $g$).} Так как
	$h(\xi)-h(x)=g(f(\xi))-g(f(x))$, а $g$ дифференцируема в точке $y=f(x)$,
	\[
	g(f(\xi))-g(f(x))=B_y\bigl(f(\xi)-f(x)\bigr)+\vphi\bigl(f(\xi)-f(x)\bigr),
	\]
	где $\vphi\colon\RR^n\to\RR^m$ удовлетворяет
	$\dfrac{\|\vphi(S)\|}{\|S\|}\to 0$ при $\|S\|\to 0$ (и $\vphi(0)=0$).

	\textit{Шаг 2 (дифференцируемость $f$).} Так как $f$ дифференцируема в точке $x$,
	\[
	f(\xi)-f(x)=A_x(\xi-x)+\psi(\xi-x),
	\]
	где $\psi\colon\RR^d\to\RR^n$ удовлетворяет
	$\dfrac{\|\psi(u)\|}{\|u\|}\to 0$ при $\|u\|\to 0$.
	Подставляя это в линейную часть из шага~1 и используя линейность $B_y$,
	\[
	B_y\bigl(f(\xi)-f(x)\bigr)
	=B_y A_x(\xi-x)+B_y\,\psi(\xi-x).
	\]
	Так как $\|B_y\,\psi(\xi-x)\|\le\|B_y\|\cdot\|\psi(\xi-x)\|$, второе слагаемое
	есть $o(|\xi-x|)$.

	\textit{Шаг 3 (оценка $\vphi$).} Из шага~2 следует, что
	$\|f(\xi)-f(x)\|\le\|A_x\|\cdot|\xi-x|+o(|\xi-x|)$, то есть
	$\|f(\xi)-f(x)\|=\bigO(|\xi-x|)$. Рассмотрим два случая при $\xi\to x$:
	\begin{itemize}
		\item если $f(\xi)=f(x)$, то $\vphi(f(\xi)-f(x))=\vphi(0)=0$;
		\item если $f(\xi)\ne f(x)$, то
		\[
		\frac{\|\vphi(f(\xi)-f(x))\|}{|\xi-x|}
		=
		\underbrace{\frac{\|\vphi(f(\xi)-f(x))\|}{\|f(\xi)-f(x)\|}}_{\to\,0}
		\cdot
		\underbrace{\frac{\|f(\xi)-f(x)\|}{|\xi-x|}}_{\text{ограничено}}\;\to\;0.
		\]
	\end{itemize}
	В обоих случаях $\vphi(f(\xi)-f(x))=o(|\xi-x|)$.

	Собирая всё вместе,
	\[
	h(\xi)-h(x)=B_y A_x(\xi-x)+o(|\xi-x|),
	\]
	что и означает дифференцируемость $h$ в точке $x$ с $d_x h=B_y A_x=d_{f(x)}g\circ d_xf$.
\end{proof}

\subsection*{Матричная форма}

В стандартных базисах дифференциалу отвечает матрица Якоби
$J_f(x)=\bigl(\partial f_i/\partial x_j(x)\bigr)$. Поскольку матрица композиции линейных
отображений есть произведение матриц,
\[
J_{g\circ f}(x)=J_g(f(x))\,J_f(x),
\]
а в координатах
\[
\frac{\partial (g_i\circ f)}{\partial x_j}(x)
=\sum_{k=1}^{n}\frac{\partial g_i}{\partial y_k}(f(x))\,\frac{\partial f_k}{\partial x_j}(x).
\]

\subsection*{Случай скалярной функции и градиент}

Пусть $f\colon\Omega\subset\RR^d\to\RR$ дифференцируема в точке $x$. Тогда
$d_xf\colon\RR^d\to\RR$ --- линейный функционал, и его действие записывается через
скалярное произведение:
\[
d_xf(h)=\langle\nabla f(x),\,h\rangle=\sum_{j=1}^{d}\frac{\partial f}{\partial x_j}(x)\,h_j,
\]
где \textbf{градиент} функции $f$ в точке $x$ --- вектор
\[
\nabla f(x)=
\left(
\frac{\partial f}{\partial x_1}(x),
\dots,
\frac{\partial f}{\partial x_d}(x)
\right).
\]
То, что $j$-я компонента $d_xf$ равна $\partial f/\partial x_j(x)$, следует из подстановки
$h=e_j$ в определение дифференцируемости (необходимое условие дифференцируемости).

\subsection*{Производная по направлению}

\begin{Definition}{Производная по направлению}{}
	Пусть $x$ --- внутренняя точка $\Omega$, $h\in\RR^d$ (часто берут $|h|=1$).
	\emph{Производной функции $f\colon\Omega\to\RR$ по направлению $h$} в точке $x$
	называется предел
	\[
	\partial_h f(x):=\lim_{t\to 0}\frac{f(x+th)-f(x)}{t},
	\]
	если он существует. Для координатного направления $h=e_j$ это в точности частная
	производная $\partial_j f(x)$.
\end{Definition}

\begin{Statement}{Вычисление производной по направлению}{}
	Если $f$ дифференцируема в точке $x$, то производная по любому направлению $h$
	существует и равна
	\[
	\partial_h f(x)=d_xf(h)=\langle\nabla f(x),\,h\rangle=\sum_{j=1}^{d}\frac{\partial f}{\partial x_j}(x)\,h_j.
	\]
\end{Statement}

\begin{proof}
	Рассмотрим путь $\gamma(t)=x+th$, $\gamma\colon(-\eps,\eps)\to\Omega$, и функцию
	одной переменной $\psi_h(t)=f(\gamma(t))=f(x+th)$. Отображение $\gamma$
	дифференцируемо, $d_t\gamma(s)=sh$, то есть $\gamma'(t)=h$. По правилу цепочки
	$\psi_h$ дифференцируема в нуле и
	\[
	\psi_h'(0)=d_xf\bigl(\gamma'(0)\bigr)=d_xf(h)=\langle\nabla f(x),h\rangle.
	\]
	С другой стороны, по определению производной функции одной переменной,
	\[
	\psi_h'(0)=\lim_{t\to0}\frac{\psi_h(t)-\psi_h(0)}{t}
	=\lim_{t\to0}\frac{f(x+th)-f(x)}{t}=\partial_h f(x).
	\]
	Сравнивая, получаем требуемое равенство.
\end{proof}

\begin{Remark}{Важность дифференцируемости}{}
	Само по себе существование производных по всем направлениям \emph{не} влечёт
	дифференцируемости и даже непрерывности: например, для
	$f(x_1,x_2)=\dfrac{x_1 x_2^2}{x_1^2+x_2^4}$ (с $f(0,0)=0$) производная по любому
	прямолинейному направлению в нуле равна $0$, однако вдоль параболы $x_1=x_2^2$
	функция стремится к $1/2$, так что $f$ даже не непрерывна в нуле. Формула
	$\partial_h f(x)=\langle\nabla f(x),h\rangle$ гарантируется именно
	дифференцируемостью.
\end{Remark}

\subsection*{Геометрический смысл градиента}

% [восстановлено] — направление наискорейшего роста и оценка через неравенство
% Коши--Буняковского--Шварца.

Пусть $f$ дифференцируема в точке $x$ и $\nabla f(x)\ne 0$. Для единичного вектора
$h$ ($|h|=1$) по неравенству Коши--Буняковского--Шварца
\[
\partial_h f(x)=\langle\nabla f(x),h\rangle\le |\nabla f(x)|\cdot|h|=|\nabla f(x)|,
\]
причём равенство достигается тогда и только тогда, когда
$h=\dfrac{\nabla f(x)}{|\nabla f(x)|}$. Аналогично минимум
$\partial_h f(x)=-|\nabla f(x)|$ достигается при противоположном направлении.

Таким образом, \textbf{градиент указывает направление наискорейшего возрастания}
функции, а его длина $|\nabla f(x)|$ равна максимальной скорости роста. В частности,
по направлениям, ортогональным $\nabla f(x)$, производная равна нулю (касательные
направления к линии уровня).
